\chapter{多电子原子的能级和光谱}

\problem{4.1}{假设 He 原子中两个电子的自旋不是 $1/2$ 而是 $1$，试定性讨论 He 原子的基态。}
% layout
\answerarea

\problem{4.2}{氦原子 $n=2$ 的激发态能级如图所示，定性解释 $\Delta E_1$ 和 $\Delta E_2$ 产生的物理原因，并比较两者的大小。
\fig{4.2}}
% layout
\answerarea
\newpage

\problem{4.3}{分别写出氖原子（$Z=10$）、镁原子（$Z=12$）、磷原子（$Z=15$）、钴原子（$Z=27$）和锗原子（$Z=32$）基态的电子组态。}
% layout
\answerarea

\problem{4.4}{计算由两个 d 电子（分等效 d 电子及非等效 d 电子两种情况）组成的电子组态能级的简并度。}
% layout
\answerarea
\newpage

\problem{4.5}{试判断下列电子组态和原子态是否可能存在并说明理由：
\begin{enumerate}
\item 电子组态：$1\mathrm p2\mathrm d$、$2\mathrm s3\mathrm f$、$2\mathrm p3\mathrm d$；
\item 原子态：$^1\mathrm P_2$、$^3\mathrm F_3$。
\end{enumerate}}
% layout
\answerarea

\problem{4.6}{按 $LS$ 耦合写出三个等效 p 电子组态可能的原子态，并计算其全部可能的状态数目。}
% layout
\answerarea
\newpage

\problem{4.7}{已知某原子的一个能级为三重结构，且随能量的增加，两个能级间隔之比为 $3:5$，试由朗德间隔定则确定这些能级的 $S$、$L$ 和 $J$ 值，并写出状态符号。}
% layout
\answerarea

\problem{4.8}{分别写出 Al（$Z=13$）、Mg（$Z=12$）、Ti（$Z=22$）基态的电子组态，以及每个原子基态的原子态。}
% layout
\answerarea
\newpage

\problem{4.9}{已知铑原子（$Z=45$）的基态电子组态为 $1\mathrm s^22\mathrm s^22\mathrm p^63\mathrm s^23\mathrm p^63\mathrm d^{10}4\mathrm s^24\mathrm p^64\mathrm d^85\mathrm s^1$。
\begin{enumerate}
\item 试按 $LS$ 耦合确定其所有可能的原子态；
\item 由洪德定则确定其基态原子态。
\end{enumerate}}
% layout
\answerarea

\problem{4.10}{铅原子基态的两个价电子都在 $6\mathrm p$ 轨道。若其中一个价电子被激发到 $7\mathrm p$ 轨道，则可组成哪些原子态？（电子间的相互作用属 $jj$ 耦合。）}
% layout
\answerarea
\newpage

\problem{4.11}{在斯特恩--盖拉赫实验中，原子态的氢从温度为 $400\ \mathrm K$ 的炉中射出，在屏上接收到两条氢线，其间距为 $0.6\ \mathrm{cm}$。若把氢原子换成氯原子，其他条件不变，那么在屏上可以接收到几条氯线？其相邻两束的间距是多少？}
% layout
\answerarea

\problem{4.12}{在磁感应强度为 $2.0\ \mathrm T$ 的磁场中，钙的 $^1\mathrm D_2\to{}^1\mathrm P_1$ 的谱线（$\lambda=732.6\ \mathrm{nm}$）分裂成相距为 $2.8\times10^{10}\ \mathrm{Hz}$ 的三个成分，试计算电子的荷质比。}
% layout
\answerarea
\newpage

\problem{4.13}{\begin{enumerate}
\item 试计算处于 P 态的原子的朗德因子 $g$；
\item 证明所有 S 态（自旋单态除外）的朗德因子 $g=2$；
\item 证明所有自旋单态的朗德因子 $g=1$。
\end{enumerate}}
% layout
\answerarea

\newpage

\problem{4.14}{指出下列原子辐射跃迁中哪些是电偶极允许跃迁，哪些是禁戒跃迁，并说明后者所违背的选择定则：
\begin{enumerate}
\item 氦 $1\mathrm s2\mathrm p\ {}^1\mathrm P_1\to1\mathrm s^2\ {}^1\mathrm S_0$；
\item 氦 $1\mathrm s2\mathrm p\ {}^3\mathrm P_1\to1\mathrm s^2\ {}^1\mathrm S_0$；
\item 碳 $3\mathrm p3\mathrm s\ {}^3\mathrm P_1\to2\mathrm p^2\ {}^3\mathrm P_0$；
\item 碳 $2\mathrm p3\mathrm s\ {}^3\mathrm P_0\to2\mathrm p^2\ {}^3\mathrm P_0$；
\item 钠 $2\mathrm p^64\mathrm d\ {}^2\mathrm D_{5/2}\to2\mathrm p^63\mathrm p\ {}^2\mathrm P_{1/2}$。
\end{enumerate}}
% layout
\answerarea
\newpage

\problem{4.15}{已知镁原子（$Z=12$）是二价原子，且符合 $LS$ 耦合。问：由价电子组态 $3\mathrm p4\mathrm p$ 直接跃迁到 $3\mathrm s3\mathrm p$ 有多少种辐射跃迁？用原子态符号表示这些可能的跃迁。}
% layout
\answerarea

\problem{4.16}{动能 $T=40\ \mathrm{keV}$ 的电子所产生的最短的 X 射线的波长为 $0.0311\ \mathrm{nm}$，试求普朗克常数。}
% layout
\answerarea
\newpage

\problem{4.17}{高速电子打在铑（$Z=45$）靶上，产生 X 连续光谱的短波限 $\lambda_{\min}=0.062\ \mathrm{nm}$，试问此时能否观察到 K 线系标识谱线？}
% layout
\answerarea

\problem{4.18}{装有钴靶的 X 射线管产生的 X 射线由较强的 Co 原子 K 线系和较弱的杂质 K 线系组成。Co 原子 $\mathrm K_\alpha$ 线波长为 $0.1791\ \mathrm{nm}$，而杂质的 $\mathrm K_\alpha$ 线波长为 $0.2291\ \mathrm{nm}$ 和 $0.1542\ \mathrm{nm}$。试利用莫塞莱定律计算两种杂质的原子序数。它们是什么元素？}
% layout
\answerarea
\newpage

